\section{The second-order theory}\label{sec:secondorder}

Theorem~\ref{thm:first} decides which face holds the mode up to factors
$N^{o(1)}$. To locate a crossing between 2 faces up to $o(1)$ in $T$ we need
the constant in front of $N^{-(k-1)}e^{-T\lambda_F}$, and in this section we
find it for every face that satisfies (H). Inside a face $F$ with $k$ states
we number the states of $F$ as $1,\dots,k$ and use
$\alpha'=(\alpha_2,\dots,\alpha_k)$ as coordinates on $\Delta_F$. Gradients
and Hessians in $\alpha$ are taken in these coordinates, $|\cdot|$ is the
Euclidean norm, $q=(q_i)_{i\in F}$, and $\mu\cdot v=\sum_{i\in F}\mu_iv_i$.

\subsection{The tilted chain and the constant}

\begin{definition}[tilt and constants]\label{def:tilt}
Let $F$ be irreducible with $k\ge2$. For
$\alpha\in\Delta_F^\circ$ let $r_\alpha$ attain the supremum defining
$I_F(\alpha)$ (Proposition~\ref{prop:rate}(b)), and put
$\theta(\alpha)=A_Fr_\alpha/r_\alpha$ and $l_\alpha=\alpha/r_\alpha$
(entrywise). The \emph{tilted generator} is
$\widehat Q_\alpha=R^{-1}(A_F-\operatorname{diag}\theta(\alpha))R$ with
$R=\operatorname{diag}r_\alpha$. With $D=\operatorname{diag}\alpha$ let
\[
Z_\alpha=(\mathbf 1\alpha^\top-\widehat Q_\alpha)^{-1}-\mathbf 1\alpha^\top,
\qquad\Gamma_\alpha=DZ_\alpha+Z_\alpha^\top D,
\]
and let $\Sigma_\alpha$ be $\Gamma_\alpha$ with the row and column of state
$1$ deleted. Finally let $K_{\{i\}}=\mu_i$ for a single state, and
\[
C_F(\alpha)=\frac{(\mu\cdot r_\alpha)(l_\alpha\cdot\mathbf 1)}
{(2\pi)^{(k-1)/2}\sqrt{\det\Sigma_\alpha}},\qquad K_F=C_F(\alpha^*_F).
\]
\end{definition}

The vector $r_\alpha$ is unique up to a positive factor, and nothing in
Definition~\ref{def:tilt} changes when it is scaled. Note that
$l_\alpha\cdot r_\alpha=1$ and $I_F(\alpha)=\langle q-\theta(\alpha),\alpha\rangle$.
Since $(A_F-\operatorname{diag}\theta(\alpha))r_\alpha=0$ with $r_\alpha>0$, and
a positive eigenvector of an irreducible Metzler matrix belongs to its Perron
root~\cite{bermanplemmons1994}, $\theta(\alpha)$ lies in the set $V_F$ defined
before Lemma~\ref{lem:upper}. The vector $r_\alpha$
minimizes $u\mapsto\sum_{i\ne j}\alpha_iq_{ij}u_j/u_i$ over $u>0$, and at
$u=r_\alpha$ the derivative in $u_j$ is
$\sum_il_{\alpha,i}q_{ij}-l_{\alpha,j}\theta_j(\alpha)$. Hence
\begin{equation}\label{eq:stationary}
l_\alpha^\top\bigl(A_F-\operatorname{diag}\theta(\alpha)\bigr)=0 .
\end{equation}
The off-diagonal entries $q_{ij}r_{\alpha,j}/r_{\alpha,i}$ of
$\widehat Q_\alpha$ have the pattern of $A_F$, and $\widehat Q_\alpha\mathbf 1=0$.
By~\eqref{eq:stationary}, $\alpha^\top\widehat Q_\alpha=0$. So
$\widehat Q_\alpha$ is an irreducible generator with stationary law $\alpha$.
If $(\mathbf 1\alpha^\top-\widehat Q_\alpha)v=0$, then $\alpha^\top v=0$ and
$\widehat Q_\alpha v=0$, so $v=0$. Thus $Z_\alpha$ exists, and
\begin{equation}\label{eq:Zprops}
Z_\alpha\mathbf 1=0,\qquad\alpha^\top Z_\alpha=0,\qquad
-\widehat Q_\alpha Z_\alpha=I-\mathbf 1\alpha^\top .
\end{equation}
So $\Gamma_\alpha\mathbf 1=0$. (The matrix $\Gamma_\alpha$ is the asymptotic
covariance, per unit time, of the occupation times of the chain with
generator $\widehat Q_\alpha$, but we do not use this.) As $\Gamma_\alpha$ is
symmetric with $\Gamma_\alpha\mathbf 1=0$, its adjugate is a multiple of
$\mathbf 1\mathbf 1^\top$, so deleting another state instead of state $1$
gives the same $\det\Sigma_\alpha$. At $\alpha^*_F$ we have
$I_F=\lambda_F$ (Proposition~\ref{prop:rate}(a)), and $u=r_F$ gives the value
$\lambda_F$ in the supremum. So $r_{\alpha^*_F}$ is a multiple of $r_F$,
$\theta(\alpha^*_F)=q-\lambda_F\mathbf 1$, and $l_{\alpha^*_F}$ is the
matching multiple of $l_F$. In the local law below, $\mu\cdot r_\alpha$
accounts for the start of the path, $l_\alpha\cdot\mathbf 1$ for its free end,
and $\Sigma_\alpha$ for the Gaussian spread of the occupation fractions.

For conductances $c_{ij}=c_{ji}$ on $F$ let
$\tree(c)=\sum_{\mathcal T}\prod_{\{i,j\}\in\mathcal T}c_{ij}$, the sum over
the spanning trees $\mathcal T$ of the complete graph on $F$.

\begin{proposition}[the constant]\label{prop:constant}
Let $F$ be irreducible with $k\ge2$.
\begin{enumerate}[(a)]
\item The maps $\alpha\mapsto\theta(\alpha)$, $I_F(\alpha)$, $\Sigma_\alpha$,
$r_\alpha$ and $l_\alpha$ (with $\sum_ir_{\alpha,i}=1$) are real-analytic on
$\Delta_F^\circ$, $\Sigma_\alpha$ is positive definite, and
$\nabla^2I_F(\alpha)=\Sigma_\alpha^{-1}$. Hence
$\det\Sigma_\alpha\cdot\det\nabla^2I_F(\alpha)=1$ and
$C_F(\alpha)=(\mu\cdot r_\alpha)(l_\alpha\cdot\mathbf 1)
\sqrt{\det\nabla^2I_F(\alpha)/(2\pi)^{k-1}}$.
\item Suppose that $\nu_iq_{ij}=\nu_jq_{ji}$ for $i,j\in F$ and some
$\nu\in(0,\infty)^F$. Then $r_\alpha$ is a multiple of $\sqrt{\alpha/\nu}$
and $l_\alpha$ of $\sqrt{\alpha\nu}$, and $\widehat Q_\alpha$ is reversible
with conductances $c_{ij}(\alpha)=\sqrt{\alpha_i\alpha_j}\,q_{ij}\sqrt{\nu_i/\nu_j}$.
Also $I_F(\alpha)=\sum_iq_i\alpha_i-\sum_{i\ne j}c_{ij}(\alpha)$, and
\[
\det\Sigma_\alpha=\frac{2^{k-1}(\prod_i\alpha_i)^2}{\tree(c(\alpha))},\qquad
C_F(\alpha)=\frac{(\mu\cdot r_\alpha)(l_\alpha\cdot\mathbf 1)}{(4\pi)^{(k-1)/2}}
\,\frac{\sqrt{\tree(c(\alpha))}}{\prod_i\alpha_i}.
\]
\end{enumerate}
\end{proposition}

\begin{proof}
(a) For $\vartheta\in\mathbb R^F$ let
$\Lambda(\vartheta)=\varrho(Q_F+\operatorname{diag}\vartheta)$, so
$\Lambda(\vartheta+c\mathbf 1)=\Lambda(\vartheta)+c$. This is a simple
eigenvalue with positive right and left eigenvectors $\hat r,\hat l$,
$\hat l\cdot\hat r=1$~\cite{bermanplemmons1994}. By the implicit function
theorem applied to the characteristic polynomial, $\Lambda$ and these vectors
are real-analytic in $\vartheta$. Put $\pi=\hat l\circ\hat r$ and
$G=\hat R^{-1}(Q_F+\operatorname{diag}\vartheta-\Lambda(\vartheta)I)\hat R$
with $\hat R=\operatorname{diag}\hat r$. As above, $G$ is an irreducible
generator with stationary law $\pi$, and we define $Z$ and $\Gamma$ from
$(G,\pi)$ as $Z_\alpha$ and $\Gamma_\alpha$ are defined from
$(\widehat Q_\alpha,\alpha)$. Then~\eqref{eq:Zprops} holds for them.

\emph{Step 1. $\nabla\Lambda(\vartheta)=\pi$ and
$\nabla^2\Lambda(\vartheta)=\Gamma$.} Since $\hat R$ commutes with diagonal
matrices, $\Lambda(\vartheta+\eta)=\Lambda(\vartheta)+\kappa(\eta)$ with
$\kappa(\eta)=\varrho(G+\operatorname{diag}\eta)$. Let $v(\eta)$ be its
eigenvector with $\pi^\top v=1$, so $v(0)=\mathbf 1$ and $\kappa(0)=0$, and
write subscripts for derivatives at $\eta=0$. Differentiating
$(G+\operatorname{diag}\eta)v=\kappa v$ in $\eta_j$ gives
$Gv_j+e_j=\kappa_j\mathbf 1$. Multiplying by $\pi^\top$ gives
$\kappa_j=\pi_j$, and then $v_j=Ze_j$ by~\eqref{eq:Zprops}, since
$\pi^\top v_j=0$. Differentiating again, in $\eta_i$, gives
$Gv_{ij}+\operatorname{diag}(e_i)v_j+\operatorname{diag}(e_j)v_i
=\kappa_{ij}\mathbf 1+\kappa_iv_j+\kappa_jv_i$, and multiplying by $\pi^\top$
gives $\kappa_{ij}=\pi_iZ_{ij}+\pi_jZ_{ji}=\Gamma_{ij}$.

\emph{Step 2. $\eta^\top\Gamma\eta>0$ unless $\eta\in\mathbb R\mathbf 1$.} Let
$u=Z\eta$ and $f=\eta-(\pi\cdot\eta)\mathbf 1$. By~\eqref{eq:Zprops},
$-Gu=f$ and $\pi\cdot u=0$, so
$\eta^\top\Gamma\eta=2\sum_i\pi_if_iu_i=-2\sum_i\pi_iu_i(Gu)_i
=\sum_{i\ne j}\pi_iG_{ij}(u_i-u_j)^2$, using that the rows of $G$ sum to
zero and $\pi^\top G=0$. The graph of $G$ is strongly connected, so this
vanishes only if $u$ is constant, hence $u=0$ and $f=0$. In particular
$\Gamma$ without the row and column of state $1$ is positive definite.

\emph{Step 3.} Let $\Lambda_0(\beta)=\Lambda(0,\beta)$ for
$\beta\in\mathbb R^{k-1}$. By Steps 1 and 2 its Hessian is positive definite
everywhere. So $\nabla\Lambda_0$ is injective, and by the inverse function
theorem it has a real-analytic inverse on its open image. Fix
$\alpha\in\Delta_F^\circ$ and put $\vartheta(\alpha)=q-\theta(\alpha)$. Then
$Q_F+\operatorname{diag}\vartheta(\alpha)=A_F-\operatorname{diag}\theta(\alpha)$
has Perron root $0$ and, by~\eqref{eq:stationary}, Perron vectors $r_\alpha$
and $l_\alpha$. So at $\vartheta(\alpha)$ we have $\pi=\alpha$,
$G=\widehat Q_\alpha$ and $\Gamma=\Gamma_\alpha$. Let $\beta(\alpha)$ be the
last $k-1$ coordinates of $\vartheta(\alpha)-\vartheta_1(\alpha)\mathbf 1$.
By the shift property, $\nabla\Lambda_0(\beta(\alpha))=\alpha'$ and
$\nabla^2\Lambda_0(\beta(\alpha))=\Sigma_\alpha$. Hence
$\beta(\alpha)=(\nabla\Lambda_0)^{-1}(\alpha')$ is real-analytic, and so are
$\vartheta_1(\alpha)=-\Lambda_0(\beta(\alpha))$, $\theta(\alpha)$, its Perron
vectors and $\Sigma_\alpha$. Moreover
$I_F(\alpha)=\langle\vartheta(\alpha),\alpha\rangle-\Lambda(\vartheta(\alpha))
=\langle\beta(\alpha),\alpha'\rangle-\Lambda_0(\beta(\alpha))$. Differentiating
and using $\nabla\Lambda_0(\beta)=\alpha'$ gives $\nabla I_F=\beta$, and then
$\nabla^2I_F=D\beta=\Sigma_\alpha^{-1}$.

(b) Put $u=\sqrt{\alpha/\nu}$ and $l=\alpha/u=\sqrt{\alpha\nu}$. Then
$\sum_il_iq_{ij}=\sum_i\sqrt{\alpha_i\nu_i}\,q_{ij}$ and
$l_j(A_Fu)_j/u_j=\nu_j\sum_iq_{ji}\sqrt{\alpha_i/\nu_i}$ agree because
$\nu_jq_{ji}=\nu_iq_{ij}$. So $u$ is a critical point of
$\sum_{i\ne j}\alpha_iq_{ij}u_j/u_i$, which is convex in $w=\log u$, and
$r_\alpha$ is a multiple of $u$. Then
$\alpha_i(\widehat Q_\alpha)_{ij}=l_iq_{ij}u_j=c_{ij}(\alpha)$, which is
symmetric by reversibility, and $\alpha_i\theta_i(\alpha)=\sum_jc_{ij}(\alpha)$
gives $I_F$, in agreement with Proposition~\ref{prop:rate}(c).

For the determinant let $\mathcal L=-D\widehat Q_\alpha$, the Laplacian of the
conductances. Fix $\eta$ with $\eta_1=0$ and let $u=Z_\alpha\eta$ and
$\zeta=D\eta-(\alpha\cdot\eta)\alpha$. By~\eqref{eq:Zprops},
$\mathcal Lu=\zeta$ and $\alpha\cdot u=0$, so
$\eta^\top\Gamma_\alpha\eta=2(D\eta)^\top u=2\zeta^\top u$. Let
$v=u-u_1\mathbf 1$. As $\zeta\cdot\mathbf 1=0$ and $\mathcal L\mathbf 1=0$,
$\zeta^\top u=\zeta^\top v$ and $\mathcal Lv=\zeta$, and since $v_1=0$ the
rows $2,\dots,k$ give $\mathcal L_{(1)}v'=\zeta'$. Here $\mathcal L_{(1)}$ is
$\mathcal L$ without row and column $1$, primes denote the coordinates
$2,\dots,k$, and $\zeta'=(D'-\alpha'\alpha'^\top)\eta'$ with
$D'=\operatorname{diag}\alpha'$. Hence
$\Sigma_\alpha=2(D'-\alpha'\alpha'^\top)\mathcal L_{(1)}^{-1}(D'-\alpha'\alpha'^\top)$.
By the matrix determinant lemma,
$\det(D'-\alpha'\alpha'^\top)=\prod_{i\ge2}\alpha_i\,(1-\sum_{i\ge2}\alpha_i)
=\prod_i\alpha_i$, and by the matrix-tree theorem~\cite{bollobas1998}
$\det\mathcal L_{(1)}=\tree(c(\alpha))$. (The weighted form follows from the
form for multigraphs, since both sides are polynomials in the $c_{ij}$.)
\end{proof}

\subsection{The local law and the face maxima}

Inside the face the occupation fractions fluctuate on the scale $T^{-1/2}$ in
$k-1$ directions. So a single composition near $\alpha$ gets the Gaussian
factor $T^{(k-1)/2}$ times the volume $N^{-(k-1)}$ of one lattice cell. This
factor gives the $\frac12$ in the crossings below.

\begin{theorem}[local law]\label{thm:local}
Let $F$ satisfy (H) with $k\ge2$, and let $\mathcal K\subset\Delta_F^\circ$ be
compact. Uniformly for $\alpha\in\mathcal K$ and $T\ge1$,
\[
f_T(\alpha)=T^{(k-1)/2}C_F(\alpha)e^{-TI_F(\alpha)}
\Bigl(1+\frac{c(\alpha)}T+O_{\mathcal K}(T^{-3/2})\Bigr),
\]
where $c(\alpha)$ is real-analytic on $\Delta_F^\circ$ and is given by
Lemma~\ref{lem:laplace} applied to the torus integral of
Lemma~\ref{lem:torus}. For $F=\{i\}$, $f_T=\mu_ie^{-q_iT}$ exactly.
\end{theorem}

The proof is in Appendix~\ref{app:torus}. We state only the first 2 terms,
which is all we use. The hypothesis $\mu(F)>0$ is needed for the relative
form, since $f_T=C_F=0$ when $\mu(F)=0$. Note that $f_T(\alpha)$ is continuous
and positive in $(T,\alpha)$. Indeed $H$ in Lemma~\ref{lem:continuum} is a
power series with coefficients $0\le W(m)\le k^Mq_{\max}^{M-1}$, so it
converges everywhere, and some $W(m)$ with every $m_i\ge1$ is positive since
$F$ is irreducible and $\mu(F)>0$.

\begin{lemma}[discrete and continuum]\label{lem:discrete}
Let $F$ satisfy (H) with $k\ge2$, and let $\mathcal K\subset\Delta_F^\circ$ be
compact. Uniformly for $\operatorname{supp}n=F$ with $n/N\in\mathcal K$ and
$1\le T\le N^{1/2}$,
$P_N(n)=N^{-(k-1)}f_T(n/N)(1+O_{\mathcal K}(T^2/N))$.
\end{lemma}

The proof is in Appendix~\ref{app:torus}.

\begin{theorem}[face maxima]\label{thm:facemax}
Let $F$ satisfy (H). If $k\ge2$, then uniformly for $1\le T\le N^{1/3}$,
\[
M_F=N^{-(k-1)}T^{(k-1)/2}K_Fe^{-T\lambda_F}\bigl(1+O(1/T+T^2/N)\bigr),
\]
and every maximizer $\hat n$ satisfies $\hat n/N=\alpha^*_F+O(1/T)$. For $k=1$,
$M_{\{i\}}=\mu_i(1-hq_i)^{N-1}=\mu_ie^{-Tq_i}(1+O(T^2/N))$. In both cases,
uniformly for $1\le T\le N^{1/3}$,
\begin{equation}\label{eq:logmax}
\log M_F=-(k-1)L+\tfrac{k-1}2\log T+\log K_F-T\lambda_F+O(1/T+T^2/N).
\end{equation}
\end{theorem}

\begin{proof}
For $k=1$ the only path is $i^N$, and
$(N-1)\log(1-hq_i)=-Tq_i+O(T^2/N)$. Let $k\ge2$, put
$X=N^{-(k-1)}T^{(k-1)/2}K_Fe^{-T\lambda_F}$ and
$x=\alpha'-\alpha^{*\prime}_F$. By Proposition~\ref{prop:constant}(a),
$I_F$ and $\log C_F$ are real-analytic on $\Delta_F^\circ$ (note that
$\mu\cdot r_\alpha>0$), $\nabla I_F(\alpha^*_F)=0$ at the interior minimum,
and $\nabla^2I_F>0$. So there are a closed ball
$U=\{|x|\le\rho_0\}\subset\Delta_F^\circ$ and constants $c_U,C_U,L',L''$ such
that on $U$ we have $c_U|x|^2\le I_F-\lambda_F\le C_U|x|^2$,
$|\log(C_F/K_F)-\gamma\cdot x|\le L'|x|^2$ with
$\gamma=\nabla\log C_F(\alpha^*_F)$, and $\log C_F$ is $L''$-Lipschitz.

\emph{Outside $U$.} As a supremum of continuous functions, $I_F$ is lower
semicontinuous on $\Delta_F$. By Proposition~\ref{prop:rate}(a) it is at least
$\lambda_F+3\eta$ on $\Delta_F\setminus U^\circ$ for some $\eta>0$. So each
$\beta\in\Delta_F\setminus U^\circ$ has some $\theta\in V_F$ with
$\langle q-\theta,\cdot\rangle>\lambda_F+2\eta$ near $\beta$, and by
compactness finitely many $\theta_1,\dots,\theta_J$ suffice. With the right
$\theta_j$ and $T^2/N\le1$, Lemma~\ref{lem:upper} gives
$P_N(n)\le C(T/N)^{k-1}e^{-T(\lambda_F+2\eta)}=X\cdot O(T^{(k-1)/2}e^{-2\eta T})$
for $\operatorname{supp}n=F$ and $n/N\notin U^\circ$.

\emph{Upper bound inside $U$.} By Theorem~\ref{thm:local} and
Lemma~\ref{lem:discrete} with $\mathcal K=U$, for $\alpha=n/N\in U$,
\begin{equation}\label{eq:insideU}
P_N(n)=X\,\tfrac{C_F(\alpha)}{K_F}\,e^{-T(I_F(\alpha)-\lambda_F)}(1+E),
\qquad|E|\le C_2(1/T+T^2/N).
\end{equation}
For $T\ge2L'/c_U$ the middle factors are at most
$\exp(\sup_x[\gamma\cdot x-\frac{c_UT}2|x|^2])=e^{|\gamma|^2/(2c_UT)}=1+O(1/T)$,
and for smaller $T$ they are bounded. So $M_F\le X(1+O(1/T+T^2/N))$.

\emph{Lower bound.} Let $n^0$ have support $F$ and
$|n^0_i-N\alpha^*_{F,i}|<1$ for all $i$, so $|x^0|\le k/N$. Then
$C_F(n^0/N)\ge K_Fe^{-kL''/N}$ and
$T(I_F(n^0/N)-\lambda_F)\le C_Uk^2T/N^2$, and~\eqref{eq:insideU} gives
$P_N(n^0)\ge X(1-O(1/T+T^2/N))$, since $1/N$ and $T/N^2$ are at most
$T^2/N$. This is useful for $T\ge T_2$ with $T_2$ large. For $T\le T_2$,
$T^2/N\to0$, so Lemma~\ref{lem:discrete} gives
$P_N(n^0)\ge\frac12N^{-(k-1)}f_T(n^0/N)\ge cX$, because $f_T$ is continuous
and positive on $[1,T_2]\times U$ and $X\le CN^{-(k-1)}$ there. So $M_F/X$ lies in $[c,C]$ for $T\le T_2$
and in $[\frac12,\frac32]$ for $T\ge T_2$, which gives the first display
and~\eqref{eq:logmax}.

\emph{The maximizer.} For $T$ bounded the claim is trivial. For $T\ge T_3$
with $T_3$ large, the bound outside $U$ is below $X/2\le P_N(n^0)$, so
$\hat\alpha=\hat n/N\in U$, and $|E|\le\frac12$. Taking logarithms in
$P_N(\hat n)\ge P_N(n^0)$ gives, with $\hat x=\hat\alpha'-\alpha^{*\prime}_F$,
\[
c_UT|\hat x|^2\le L''|\hat x|+C_4(1/T+T^2/N)\le L''|\hat x|+2C_4/T,
\]
because $T\le N^{1/3}$ gives $T^2/N\le1/T$. This step uses the exponent
$\frac13$. If $L''|\hat x|\ge2C_4/T$ then $|\hat x|\le2L''/(c_UT)$,
and otherwise $|\hat x|<2C_4/(L''T)$.
\end{proof}

\subsection{Crossings and tie windows}

\begin{theorem}[the crossing equation]\label{thm:crossing}
Let $F$ and $G$ satisfy (H), with $\Delta=k_F-k_G\ge1$ and
$\delta=\lambda_G-\lambda_F>0$, and let $D_N(T)=\log M_F-\log M_G$.
\begin{enumerate}[(i)]
\item For large $N$, $D_N$ has a zero in
$[\frac{\Delta}{2\delta}L,\frac{2\Delta}{\delta}L]$.
\item For large $N$, every zero $T_N\in[1,N^{1/3}]$ of $D_N$ satisfies
\begin{equation}\label{eq:crossing}
\delta T_N+\frac\Delta2\log T_N=\Delta L-\log\frac{K_F}{K_G}
+O\Bigl(\frac1L\Bigr).
\end{equation}
\item All such zeros lie in one interval of length $O(1/L)$, and
\[
T_N=\frac\Delta\delta\Bigl[L-\frac12\log L+b_{FG}\Bigr]
+O\Bigl(\frac{\log L}L\Bigr),\qquad
b_{FG}=-\frac12\log\frac\Delta\delta-\frac1\Delta\log\frac{K_F}{K_G}.
\]
\end{enumerate}
\end{theorem}

\begin{proof}
By~\eqref{eq:logmax}, uniformly on $[1,N^{1/3}]$,
$D_N(T)=-\Delta L+\frac\Delta2\log T+\delta T+\log\frac{K_F}{K_G}+O(1/T+T^2/N)$.

(i) At $T=\frac{\Delta}{2\delta}L$ this is $-\frac\Delta2L+O(\log L)<0$, and at
$T=\frac{2\Delta}\delta L$ it is $\Delta L+O(\log L)>0$, for large $N$. For
fixed $N$, $M_F$ is a maximum of finitely many polynomials in $h=T/N$, and it
is positive because $F$ is admissible. So $D_N$ is continuous and has a zero
in between.

(ii) Let $\psi(T)=\delta T+\frac\Delta2\log T$, which is increasing. At a zero
$T_N\in[1,N^{1/3}]$, $\psi(T_N)=\Delta L-\log(K_F/K_G)+O(1)$. This is larger
than $\psi(\frac\Delta{2\delta}L)=\frac\Delta2L+O(\log L)$ for large $N$, and
$\delta T_N\le\psi(T_N)$. So $\frac\Delta{2\delta}L<T_N\le\frac{2\Delta}\delta L$,
and then $O(1/T_N+T_N^2/N)=O(1/L)$.

(iii) Let $\psi(T^\circ)=\Delta L-\log(K_F/K_G)$. As $\psi'\ge\delta$, every
zero is within $O(1/L)$ of $T^\circ$. First
$T^\circ=\frac\Delta\delta L+O(\log L)$, so
$\log T^\circ=\log L+\log\frac\Delta\delta+O(\log L/L)$, and substituting
this into $\psi(T^\circ)$ gives the explicit form.
\end{proof}

The error of the explicit form comes only from the last expansion, and it is
large at every $N$ we can compute. So computations should solve the implicit
equation~\eqref{eq:crossing}. Take the $4$-cycle with unit rates and uniform
start, $F$ an arc of $3$ states and $G$ a vertex, so that $\Delta=2$,
$\delta=\sqrt2$ and $\log(K_F/K_G)=0.964591$ (Remark~\ref{rem:C4}).
Numerically, at $N=10^{14}$ the exact crossing is $T_N=42.2633$, with
$\Delta L-\frac\Delta2\log T_N-\delta T_N=0.959070$. The root of the implicit
equation is within $0.004$ of $T_N$, while the explicit form gives
$42.2059$. We do not prove that the zero of $D_N$ is unique in general.
For their faces Papers~A and~B prove uniqueness, and
Corollaries~\ref{cor:paperA} and~\ref{cor:paperB} use that.

\begin{theorem}[tie windows]\label{thm:window}
Fix $\tau^*>0$. Let $m^*$ be the least value of $\Phi_{\tau^*}$ over the
admissible faces and $\mathcal F^*$ the set of minimizers. Assume that every
face in $\mathcal F^*$ is irreducible. By Theorem~\ref{thm:hull}(c) this holds
if $Q$ has symmetric support or $\mu$ has full support. Put $T=\tau^*(L-\frac12\log L+t)$, and for
$F\in\mathcal F^*$ let
$\ell_F(t)=\log K_F+\frac{k_F-1}2\log\tau^*-\tau^*\lambda_Ft$ and
$E(t)=\max_{G\in\mathcal F^*}\ell_G(t)$.
\begin{enumerate}[(i)]
\item Uniformly for $t$ in compact sets, for $F\in\mathcal F^*$,
$\log M_F=-m^*L+\frac{m^*}2\log L+\ell_F(t)+O(\log L/L)$.
\item There are $\varepsilon_1,\varepsilon_0>0$ such that for large $N$ and
$|T/L-\tau^*|\le\varepsilon_1$, every admissible $F\notin\mathcal F^*$ has
$M_F\le N^{-\varepsilon_0}\max_{G\in\mathcal F^*}M_G$.
\item There are $\varepsilon,K>0$ such that for large $N$ and
$K\le|t|\le\varepsilon L$, every global mode lies on a face
$F\in\mathcal F^*$ with $\ell_F=E$ on $[K,\infty)$ if $t>0$, or on
$(-\infty,-K]$ if $t<0$. Such a face has the least $\lambda_F$ in
$\mathcal F^*$ if $t>0$ and the largest if $t<0$.
\end{enumerate}
Consequently, for every $\eta>0$ and large $N$, if $|t|\le\varepsilon L$ is at
distance at least $\eta$ from every kink of $E$, then every global mode lies
on a face whose line equals $E$ near $t$. So the line of the face that
carries the global mode changes only at
$T=\tau^*(L-\frac12\log L+t_j)+o(1)$, where $t_j$ runs over the kinks of $E$,
and it changes near each kink. Faces with identical lines (the same $k$,
$\lambda$ and $K$) tie at this order and give no crossing.
\end{theorem}

\begin{proof}
Every $F\in\mathcal F^*$ is admissible and irreducible, so $\mu(F)>0$ by
Lemma~\ref{lem:admissible}(b), and $F$ satisfies (H). Write
$T=\tau^*L(1+y)$ with $y=(t-\frac12\log L)/L$, so that $|y|\le\frac12$ for
$|t|\le L/4$ and large $N$. By~\eqref{eq:logmax} and
$\tau^*\lambda_F+k_F-1=m^*$, uniformly for $|t|\le L/4$,
\begin{equation}\label{eq:windowline}
\log M_F=-m^*L+\tfrac{m^*}2\log L+\ell_F(t)+\tfrac{k_F-1}2\log(1+y)+O(1/L)
\qquad(F\in\mathcal F^*).
\end{equation}
For $t$ in a compact set, $\log(1+y)=O(\log L/L)$, which gives (i).

(ii) Let $g>0$ be the least value of $\Phi_{\tau^*}(F)-m^*$ over the admissible
$F\notin\mathcal F^*$. As $0\le\lambda_F\le q_{\max}$, for
$|\tau-\tau^*|\le\varepsilon_1=g/(4q_{\max}+4)$ we have
$\Phi_\tau(F)-\Phi_\tau(G)\ge g/2$ for such $F$ and all $G\in\mathcal F^*$.
With $\tau=T/L$, Theorem~\ref{thm:first}(ii) gives
$\log M_F\le-L\Phi_\tau(F)+(k_F-1)\log T+C$ and $\log M_G\ge-L\Phi_\tau(G)-C$.
So $\log M_F-\log M_G\le-\frac g2L+O(\log L)\le-\frac g4L$.

(iii) Take $K\ge1$ larger than every $|t_j|$, and
$\varepsilon\le\min(\frac14,\varepsilon_1/(2\tau^*))$, so that (ii) applies.
Let $t\ge K$, let $\ell_F=E$ on $[K,\infty)$, and let $G\in\mathcal F^*$ with
$\ell_G\ne\ell_F$. The line with the least $\lambda$ has the largest slope, so
$\lambda_F\le\lambda_G$. By~\eqref{eq:windowline},
$\log M_F-\log M_G=\ell_F(t)-\ell_G(t)+\frac{k_F-k_G}2\log(1+y)+O(1/L)$. If
$\lambda_G=\lambda_F$, then $k_G=k_F$ because both faces are tied, and the
difference is $\log(K_F/K_G)+O(1/L)$, where $K_F>K_G$ because the 2 lines
are parallel and $\ell_G\le\ell_F$. Otherwise
$\Delta=k_F-k_G=\tau^*(\lambda_G-\lambda_F)>0$, again because both faces are
tied, and $\ell_F(t)-\ell_G(t)=\Delta t+\ell_F(0)-\ell_G(0)$. As
$|\log(1+y)|\le2|y|\le t/2$ for large $N$, the difference is at least
$\frac\Delta2t-|\ell_F(0)-\ell_G(0)|-O(1/L)>0$ if $K$ is large. The case
$t\le-K$ is the same with the order of the $\lambda$'s reversed.

For the consequence, fix $\eta>0$ and let
$S_\eta=[-K,K]\setminus\bigcup_j(t_j-\eta,t_j+\eta)$. For $G\in\mathcal F^*$,
$E-\ell_G$ is convex, piecewise linear and nonnegative, so its zero set is an
interval with kinks of $E$ as endpoints, or a single kink. So on each
component of $S_\eta$ either $\ell_G=E$ or $E-\ell_G\ge g_\eta$ for some
$g_\eta>0$. By (i) and (ii), for large $N$ and $t\in S_\eta$ a global mode
lies on a face with $\ell_F(t)=E(t)$, and for $K\le|t|\le\varepsilon L$ we use
(iii). The line equal to $E$ is constant on each component of
$[-\varepsilon L,\varepsilon L]\setminus\bigcup_j(t_j-\eta,t_j+\eta)$ and
changes at each kink. Since $\eta$ is arbitrary, the claim follows.
\end{proof}

\begin{remark}[scope of the $\frac12$]\label{rem:half}
The coefficient $\frac12$ of $\log\log N$ in Theorems~\ref{thm:crossing}
and~\ref{thm:window} is due to the power $T^{(k-1)/2}$ in
Theorem~\ref{thm:facemax}. It is the same for every pair of faces that
satisfy (H), whatever the graph, the rates, the sizes and the start, and it
needs no reversibility or aperiodicity. At a kink of $E$ where the line of
$G$ gives way to the line of $F$ as $t$ increases, we have
$\lambda_F<\lambda_G$ and $\Delta=\tau^*\delta$, and solving
$\ell_F=\ell_G$ gives $t_j=b_{FG}$. So when every tied face is irreducible,
every change of the line of the face that carries the global mode in a tie
window is at $T=\frac\Delta\delta(L-\frac12\log L+b_{FG})+o(1)$, a crossing of
Theorem~\ref{thm:crossing}. Without this assumption the $\frac12$ can fail, and
Proposition~\ref{prop:threecycle} gives the coefficients $1$ and $0$.
\end{remark}

\subsection{Papers A and B}\label{sec:papersAB}

We recall the notation of Paper~A~\cite{paperA}. For positive integers $a,b$
let $S_{a,b}(u)=\sum_wu^{s(w)}$, the sum over the words $w$ with $a$ letters
$R$ and $b$ letters $L$, where $s(w)$ is the number of switches of $w$. Paper~A
shows that $S_{a,b}(u)=P_{a+b}(a)/P_{a+b}(0)$ for its walk with $u=(1-p)/p$,
and writes it as an explicit polynomial~\cite[Eq.~(9)]{paperA}. Let
$S_N=S_{\lfloor N/2\rfloor,\lceil N/2\rceil}$, let $u_N$ be the positive root
of $S_N(u_N)=1$, and let $p_N=1/(1+u_N)$ be the unique crossing of the middle
and the endpoints~\cite[Theorem~5.1]{paperA}. Put $c_0=\frac12\log(\pi/8)$.

\begin{corollary}[Paper A]\label{cor:paperA}
Let $d=2$, $q_{12}=q_{21}=1$ and $\mu=(\frac12,\frac12)$, the walk of Paper~A
with $h=1-p$ and $T=N(1-p)$. Then $K_{\{1\}}=K_{\{2\}}=\frac12$ with
$\lambda=1$, and $K_{\{1,2\}}=\sqrt{2/\pi}$ with $\lambda=0$ and
$\det\Sigma=\frac14$. For every $N\ge2$, $M_{\{1,2\}}=M_{\{1\}}$ holds exactly
at $T=N(1-p_N)$, and for large $N$
\[
T+\tfrac12\log T=L+c_0+O(1/L)\qquad\text{at }T=N(1-p_N).
\]
So $b_{FG}=c_0$, the constant of~\cite[Theorem~5.2]{paperA}.
\end{corollary}

\begin{proof}
The full face is reversible with $\nu=\mathbf 1$, $\alpha^*=(\frac12,\frac12)$,
$r_F=\mathbf 1$ and $l_F=\frac12\mathbf 1$, so $(\mu\cdot r)(l\cdot\mathbf 1)=1$.
By Proposition~\ref{prop:constant}(b), $\tree(c)=c_{12}=\frac12$ and
$\det\Sigma=2\cdot\frac1{16}\cdot2=\frac14$, so
$K_{\{1,2\}}=2/\sqrt{2\pi}$. By~\cite[Prop.~6.1]{paperA} and the symmetry
$S_{a,b}=S_{b,a}$, every $S_{a,N-a}$ is dominated coefficientwise by $S_N$, so
the middle bin is the largest bin with support $\{1,2\}$ for every $p$. Hence
$M_{\{1,2\}}=M_{\{1\}}S_N(u)$ with $u=h/(1-h)$, and the crossing is at $p_N$.
By~\cite[Theorem~5.1]{paperA}, $N(1-p_N)\sim L$. Theorem~\ref{thm:crossing}
with $\Delta=\delta=1$ and $\log(K_{\{1,2\}}/K_{\{1\}})=\log(2\sqrt{2/\pi})=-c_0$
gives the equation.
\end{proof}

Theorem~\ref{thm:crossing}(iii) then gives
$N(1-p_N)=L-\frac12\log L+c_0+O(\log L/L)$, which is~\cite[Eq.~(7)]{paperA}.
With $c_{\{1,2\}}-c_{\{1\}}=-\frac18$, Proposition~\ref{prop:thirdorder}(ii)
turns the equation into $T+\frac12\log T=L+c_0+\frac1{8T}+O(T^{-3/2})$. This
is the equation behind~\cite[Eq.~(8)]{paperA}, and solving it recovers that
expansion, with the weaker error $O(L^{-3/2})$ in place of the
$O((\log L)^2/L^2)$ of Paper~A.

\begin{corollary}[Paper B]\label{cor:paperB}
Let $Q=J-dI$, where $J$ is the all-ones matrix, and let $\mu$ be uniform.
This is the model of Paper~B~\cite{paperB} with $r=h$ and $T=Nh$. For $|F|=k$,
$\lambda_F=d-k$, $\alpha^*_F$ is uniform on $F$,
$\det\Sigma_{\alpha^*_F}=2^{k-1}k^{1-2k}$ and
$K_F=\frac1dk^{k+1/2}(4\pi)^{-(k-1)/2}$. The crossing of a face of $k$ states
with a vertex has $\Delta=\delta=k-1$ and
$b_{FG}=a_k=\frac12\log(4\pi)-\frac{k+1/2}{k-1}\log k$, the constant
of~\cite[Theorem~3.1]{paperB}. At $\alpha\in\Delta_F^\circ$, with
$A=\sum_{i\in F}\sqrt{\alpha_i}$,
\[
I_F(\alpha)=d-A^2,\qquad
d\,C_F(\alpha)=A^{k/2+1}\Bigl(\prod_{i\in F}\alpha_i\Bigr)^{-3/4}(4\pi)^{-(k-1)/2},
\]
which is the constant $D(\alpha)$ of~\cite[Theorem~3.5]{paperB}. At
$\tau^*=1$ every face ties, and Theorem~\ref{thm:window} gives the window
of~\cite[Theorem~3.2]{paperB}.
\end{corollary}

\begin{proof}
All rates are $1$ and $q_i=d-1$, so every face is reversible with
$\nu=\mathbf 1$. By Proposition~\ref{prop:constant}(b), $r_\alpha$ and
$l_\alpha$ are multiples of $\sqrt\alpha$, so
$(\mu\cdot r_\alpha)(l_\alpha\cdot\mathbf 1)=A^2/d$, and
$c_{ij}(\alpha)=x_ix_j$ with $x_i=\sqrt{\alpha_i}$. So
$I_F(\alpha)=d-1-\sum_{i\ne j}x_ix_j=d-A^2$. For $c_{ij}=x_ix_j$ and
$S=\sum_ix_i$, the Laplacian grounded at state $1$ is
$\operatorname{diag}(x_iS)_{i\ge2}-x'x'^\top$, and the matrix determinant
lemma gives the weighted Cayley formula
$\tree(c)=S^{k-1}\prod_{i\ge2}x_i\,(1-\sum_{i\ge2}x_i/S)=x_1\cdots x_kS^{k-2}$.
So $\sqrt{\tree(c(\alpha))}=(\prod_i\alpha_i)^{1/4}A^{(k-2)/2}$, which gives
$d\,C_F(\alpha)$. By Cauchy--Schwarz, $A^2\le k$ with equality only at the
uniform $\alpha$. So $\alpha^*_F$ is uniform and $\lambda_F=d-k$, and $A=\sqrt k$
and $\tree(c)=1/k$ there give $K_F$ and $\det\Sigma$. A vertex has $K=1/d$ and
$\lambda=d-1$, so $b_{FG}=-\frac1{k-1}\log(dK_F)=a_k$.

Since $p=1-(d-1)h\ge\frac12\ge\frac1d$, the balanced bins are the largest
bins on a face~\cite[Theorem~3.6]{paperB}, so $M_F$ is Paper~B's balanced face
bin and $M_{\{i\}}$ its vertex bin. Their crossing is unique
by~\cite[Theorem~3.1]{paperB}, and there Theorem~\ref{thm:crossing} gives
$T+\frac12\log T=L+a_k+O(1/L)$, in agreement with Paper~B, whose variable is
$Nu=T/p=T+O(L^2/N)$. By Theorem~\ref{thm:local} and
Lemma~\ref{lem:discrete}, and since the vertex bin is
$\frac1de^{-(d-1)T}(1+O(T^2/N))$, the ratio of an interior bin to a vertex
bin is $d\,C_F(\alpha)T^{(k-1)/2}e^{(A^2-1)T}N^{-(k-1)}(1+O(1/T+T^2/N))$, as
in~\cite[Theorem~3.5]{paperB}. At $\tau^*=1$, $\Phi_1(F)=d-1$ for every face,
and $Q$ has symmetric support. Theorem~\ref{thm:window}(i) gives
$\log(M_F/M_{\{i\}})=\ell_F(t)-\ell_{\{i\}}(t)+O(\log L/L)=(k-1)(t-a_k)
+O(\log L/L)$, which is~\cite[Theorem~3.2]{paperB}, since the 2 ways of
writing $T$ in the window differ by $o(1)$ in $t$.
\end{proof}

So Paper~B's factor $A^{k/2+1}(\prod_i\alpha_i)^{-3/4}$ equals
$A^2\sqrt{\tree(c(\alpha))}/\prod_i\alpha_i$. The part that depends on the
graph is the square root of a weighted count of spanning trees.

\subsection{Exact examples}

\begin{proposition}[an exact pair identity]\label{prop:pair}
Let $i\ne j$ with $q_{ij}=q_{ji}=a$, $q_i=q_j=q$ and $\mu_i=\mu_j>0$. For
every $N\ge2$, $1\le m\le N-1$ and $h>0$ such that $I+hQ$ is stochastic and
$qh<1$,
\[
\frac{P_N(me_i+(N-m)e_j)}{P_N(Ne_i)}=S_{m,N-m}(u),\qquad u=\frac{ah}{1-qh}.
\]
\end{proposition}

\begin{proof}
A path with this composition is a word in $i,j$ with $m$ letters $i$. Each
stay has probability $1-qh$ and each switch $ah$, so a word with $s$ switches
has probability $\mu_i(1-qh)^{N-1-s}(ah)^s=\mu_i(1-qh)^{N-1}u^s$, while
$P_N(Ne_i)=\mu_i(1-qh)^{N-1}$. Sum over the words.
\end{proof}

The reason is that each switch replaces a factor $1-qh$ by $ah$, and nothing
else about the graph enters. So one state against a symmetric pair is
Paper~A's crossing at every $N$. We use this for $K_3$ and the cycles in
Section~\ref{sec:examples} and for sticky kernels in
Proposition~\ref{prop:sticky-pair}.

The next example shows that the irreducibility assumption cannot simply be
dropped from Theorem~\ref{thm:window}. The best composition of a reducible
face can visit a state only once. Then there is no Gaussian window in that
direction, and the power of $T$ is $1$ rather than $\frac12$. The assumption
fails at both tie points below.

\begin{proposition}[a directed $3$-cycle]\label{prop:threecycle}
Let $d=3$, $q_{12}=b>2$, $q_{23}=q_{31}=1$, all other rates $0$, and
$\mu=\delta_1$.
\begin{enumerate}[(i)]
\item The admissible faces are $\{1\}$, $\{1,2\}$ (reducible, $\lambda=1$) and
$\{1,2,3\}$ (irreducible, $\lambda=0$). Exactly
$M_{\{1\}}=(1-bh)^{N-1}$ and $M_{\{1,2\}}=bh(1-h)^{N-2}$, and the latter is
attained only at $n=(1,N-1,0)$. For the full face,
$\alpha^*=(1,b,b)/(1+2b)$, $\det\nabla^2I=(1+2b)^4/(3b^2)$ and
$K_{\{1,2,3\}}=(1+2b)^2/(2\sqrt3\,\pi b)$.
\item $M_{\{1\}}=M_{\{1,2\}}$ has exactly one root $T^{(1)}_N$ with $bh<1$. It
satisfies $(b-1)T+\log T=L-\log b+O(T^2/N)$, so
$T^{(1)}_N=\frac1{b-1}[L-\log L+\log\frac{b-1}b]+O(\log L/L)$, and the
coefficient of $\log\log N$ in the bracket is $1$.
\item Every root $T\in[L/2,2L]$ of $M_{\{1,2\}}=M_{\{1,2,3\}}$ satisfies
$T=L+\log\frac{2\sqrt3\,\pi b^2}{(1+2b)^2}+O(1/L)$, so the coefficient of
$\log\log N$ is $0$. These roots lie in an interval of length $O(1/L)$.
\item Both are changes of the global mode. For large $N$ and $T\le N^{1/3}$,
the mode is $(N,0,0)$ for $T<T^{(1)}_N$, it is $(1,N-1,0)$ between
$T^{(1)}_N$ and the smallest root in (iii), and it lies on the full face above
the largest root in (iii).
\end{enumerate}
\end{proposition}

\begin{proof}
(i) The chain moves only $1\to2\to3\to1$ and starts at $1$, so the support of
a path is $\{1\}$, $\{1,2\}$ or $\{1,2,3\}$. The components of $\{1,2\}$ are
$\{1\}$ and $\{2\}$ with exit rates $b$ and $1$. A path with support
$\{1,2\}$ is $1^{n_1}2^{n_2}$, with probability
$(1-bh)^{n_1-1}bh(1-h)^{n_2-1}$, and raising $n_1$ by one multiplies this by
$(1-bh)/(1-h)<1$. For the full face, $\lambda=0$, $r=\mathbf 1$, and the
stationary law solves $\pi_2=b\pi_1=\pi_3$. So $\alpha^*=l=(1,b,b)/s$ with
$s=1+2b$, and $(\mu\cdot r)(l\cdot\mathbf 1)=1$. The ratios $u_2/u_1$,
$u_3/u_2$, $u_1/u_3$ range over the positive triples with product $1$, so by
the inequality of arithmetic and geometric means
$I(\alpha)=\sum_iq_i\alpha_i-3(b\alpha_1\alpha_2\alpha_3)^{1/3}$. Let
$g=(\alpha_1\alpha_2\alpha_3)^{1/3}$ and $w=1/\alpha$. The Hessian of $g$ in
all 3 coordinates is $-g\,M$ with
$M=\frac13\operatorname{diag}(w)^2-\frac19ww^\top$, so in the coordinates
$(\alpha_2,\alpha_3)$, $\nabla^2I=3b^{1/3}g\,J^\top MJ$, where $J$ is the
Jacobian of $(\alpha_2,\alpha_3)\mapsto(1-\alpha_2-\alpha_3,\alpha_2,\alpha_3)$.
At $\alpha^*$, $3b^{1/3}g=3b/s$ and $w=(s,s/b,s/b)$. With $y=1/b$,
$J^\top MJ$ has diagonal entries $\frac29s^2(1+y+y^2)$ and off-diagonal
entries $\frac19s^2(2+2y-y^2)$, so
\[
\det J^\top MJ=\tfrac{s^4}{81}\bigl[4(1+y+y^2)^2-(2+2y-y^2)^2\bigr]
=\tfrac{s^4}{81}\,3y^2(2+y)^2=\tfrac{s^6}{27b^4}.
\]
Hence $\det\nabla^2I=(3b/s)^2s^6/(27b^4)=s^4/(3b^2)$, and
Proposition~\ref{prop:constant}(a) gives $K_{\{1,2,3\}}$.

(ii) Let $D_1=\log M_{\{1,2\}}-\log M_{\{1\}}
=\log(bh)+(N-2)\log(1-h)-(N-1)\log(1-bh)$. Then
$D_1'(h)=\frac1h-\frac{N-2}{1-h}+\frac{b(N-1)}{1-bh}>0$ for $bh<1$, while
$D_1\to-\infty$ as $h\to0$ and $D_1\to\infty$ as $bh\to1$. So there is
exactly one root. Expanding,
$D_1=\log b+\log T-L+(b-1)T+O(T^2/N)$, which is positive at $T=L/(b-1)$ and
negative at $T=L/(2(b-1))$. So the root lies between them, and substituting
$\log T=\log L-\log(b-1)+O(\log L/L)$ gives the expansion.

(iii) Here $\log M_{\{1,2\}}=\log b+\log T-L-T+O(T^2/N)$, and
Theorem~\ref{thm:facemax} for the full face ($k=3$, $\lambda=0$, $\mu(F)=1$)
gives $\log M_{\{1,2,3\}}=-2L+\log T+\log K_{\{1,2,3\}}+O(1/T+T^2/N)$. So
$\log M_{\{1,2\}}-\log M_{\{1,2,3\}}=L-T+\log b-\log K_{\{1,2,3\}}+O(1/T+T^2/N)$,
and every zero in $[L/2,2L]$ is within $O(1/L)$ of
$L+\log b-\log K_{\{1,2,3\}}$.

(iv) Only the 3 faces in (i) are admissible. Lemma~\ref{lem:upper} with
$\theta=q$ gives $M_{\{1,2,3\}}\le C(T/N)^2$ for $T\le N^{1/3}$. For
$T<T^{(1)}_N$, $M_{\{1\}}>M_{\{1,2\}}$ by the monotonicity of $D_1$, and
$M_{\{1\}}\ge e^{-bT-b^2T^2/N}\ge N^{-b/(b-1)-o(1)}$ since
$T^{(1)}_N\le L/(b-1)$. This beats $C(T/N)^2$ because $b/(b-1)<2$. For
$T>T^{(1)}_N$, $M_{\{1,2\}}>M_{\{1\}}$, and we compare $\{1,2\}$ with the full
face. For $T\le L/2$,
$\log M_{\{1,2\}}-\log M_{\{1,2,3\}}\ge L-T-\log T-C>0$. On $[L/2,2L]$ the
sign is given by the formula in (iii), which is positive at $L/2$ and
negative at $2L$. For $2L\le T\le N^{1/3}$ the same formula is at most
$-L+C<0$.
\end{proof}

\subsection{The next term}

The face maxima also have a $1/T$ term. It collects the $1/T$ term of the
local law at $\alpha^*_F$ and the gain from moving the maximizer by
$\Sigma\gamma/T$ toward larger $C_F$.

\begin{proposition}[the next term]\label{prop:thirdorder}
Let $F$ satisfy (H).
\begin{enumerate}[(i)]
\item The limit
$c_F=\lim_{T\to\infty}T\bigl[\sup_{\alpha\in\Delta_F^\circ}f_T(\alpha)/
(T^{(k-1)/2}K_Fe^{-T\lambda_F})-1\bigr]$ exists. For $k\ge2$,
$c_F=c(\alpha^*_F)+\frac12\gamma^\top\Sigma\gamma$, where $c$ is the
coefficient in Theorem~\ref{thm:local}, $\gamma=\nabla\log C_F(\alpha^*_F)$
and $\Sigma=\Sigma_{\alpha^*_F}$, in the coordinates $\alpha'$.
\item Let $G$ also satisfy (H), with $\Delta$ and $\delta$ as in
Theorem~\ref{thm:crossing}. At every zero $T_N\in[1,N^{1/3}]$ of
$\log M_F-\log M_G$,
\[
T_N\Bigl(\Delta L-\frac\Delta2\log T_N-\delta T_N-\log\frac{K_F}{K_G}\Bigr)
=c_F-c_G+O\bigl(T_N^{-1/2}+T_N^3/N\bigr).
\]
\item $c_F=0$ if $F$ is a single state, $c_F=-1/(8q_{ij})$ if $F=\{i,j\}$
with $q_{ij}=q_{ji}$, $q_i=q_j$ and $\mu_i=\mu_j$ (so $c_F=-\frac18$ for unit
rates), and $c_F=-(k-1)/8$ if $Q=J-dI$ and $\mu$ is constant on $F$.
\end{enumerate}
\end{proposition}

The proof is in Appendix~\ref{app:laplace}. For the arc of $3$ states and the
whole $4$-cycle, the formula in (i) is evaluated in Remark~\ref{rem:C4}.
